% §7 Discussion and Limitations — Target 1pp

\section{Discussion and Limitations}
\label{sec:discussion}

\subsection{Limitations}
\label{sec:disc:limitations}

\paragraph{Kerckhoffs scope.}
Our threat model (Section~\ref{sec:bg:lifecycle}) assumes the
attacker knows the vault's axis values, including public parameters
such as the recovery address. We do not rely on hiding attacks
(taproot leaf concealment, address non-disclosure, closed-source
construction) as security mechanisms: a vault whose security depends
on the attacker's ignorance of its construction is not
cryptographically secure by Kerckhoffs's criterion. Any result or
recommendation in this paper that is read outside this threat model
should be reinterpreted accordingly.

\paragraph{Regtest vs.\ mainnet.}
All measurements are taken on Bitcoin regtest with a single mempool.
The class \classfee{} rationality block
(Appendix~\ref{app:rationality}) scopes the fee-pinning feasibility
claim: in practice it depends on a network-layer propagation race
between the attacker's 25 descendants and the watchtower's CPFP
replacement. Regtest cannot model this race. Structural vsize
measurements are unaffected (they are identical on regtest and
mainnet), but the \emph{feasibility} of a race-sensitive attack
depends on deployment-environment parameters beyond our harness.
Mainnet-equivalent measurement on Inquisition signet is future work.

\paragraph{Bounded model checking.}
Alloy's finite-scope analysis (41 assertions across five models)
provides refutation for broad claims but does not constitute a full
proof. Mechanised proofs (Coq, Lean) of
Proposition~\ref{thm:griefing_safety} and of the feature-vulnerability
implications D1--D4 (Section~\ref{sec:impossibility}) remain open.

\paragraph{Unmerged BIPs.}
BIP-119, BIP-345, BIP-443, BIP-347, and BIP-348 are all under
review, and specifications may evolve. We tracked specifications via
pinned git hashes (see anonymised artifact); readers should verify
against the current BIP text.

\paragraph{Class \classhot{} (TM3) race dynamics are network-layer,
not Script-layer.}
Our design-space decomposition does not model the race between the
attacker's unvault-plus-withdrawal sequence and the watchtower's
recovery broadcast. Formalising $\Pr_D[\text{attacker wins race}]$
requires a separate probability space over block-arrival times,
mempool propagation delays, and detection latency---parameters that
depend on the deployment environment rather than on BIP script
semantics. Existing state-channel watchtower formalisations treat
confirmation as atomic or assume bounded-delay synchrony; none
provides off-the-shelf machinery for a pinning-composed race against
a Bitcoin mempool. We therefore treat class \classhot{} (TM3) as
uniformly applicable to all four Script extensions at the script
layer, with \catcsfs{}'s dual-binding on the hot leaf
(\axwdbind{} = \texttt{signature-dual-bind}) as the one design-level
mitigation (reducing the worst-case outcome from theft to griefing).

\paragraph{Batching operational complexity.}
The measured $r^\dagger$ gains of
Section~\ref{sec:empirical:batching} characterise the best-case
defender strategy. Real deployments face buffering (CSV-window
deadlines per pending UTXO), single-input-poisons-batch failure
modes, and, for BIP-345 authorised recovery, an online
\texttt{recoveryauth} key that enlarges the watchtower's own attack
surface. These are upper bounds on achievable safety margin, not
refutations of the thresholds.

\subsection{Open Problems}
\label{sec:disc:open}

\begin{enumerate}
  \item Mainnet-equivalent measurement on Inquisition signet.
  \item Mechanised proofs (Coq/Lean) of
    Proposition~\ref{thm:griefing_safety} and of the
    feature-vulnerability implications D1--D4
    (Section~\ref{sec:impossibility}).
  \item Impact of TRUC/BIP-431 relay policy on class \classfee{}
    (fee-channel pinning). Under TRUC the 25-descendant limit
    collapses to 1, eliminating the pinning surface at the
    relay-policy level (see Section~\ref{sec:imposs:class_fee}
    counter-factual and Appendix~\ref{app:rationality}).
  \item Recursive covenant design space beyond \opccv{}'s taptree
    preservation.
  \item Cross-chain generalisation: Ethereum account abstraction as
    an L1 analogue of the design-space decomposition.
\end{enumerate}

\providecommand{\todo}[2][]{\textbf{[TODO: #2]}}
